\subsection{RTL to Mapped Synthesis: the ACE-2 Accelerator}
\label{sec:silicon-vertical}

The campaigns above produce software, proofs, and manuscripts. ACE-2 asks whether
the same runtime holds when the artifact is hardware, where an unverified claim
survives exactly as long as it takes someone to run the tool that contradicts it.
Prior work on language models for hardware has concentrated on generating and
evaluating RTL at the module level~\citep{liu2023verilogeval} and on adapting models
to chip-design corpora~\citep{liu2023chipnemo}; the question here is different,
namely whether a runtime can carry an accelerator from model structure through
functional closure and mapped synthesis as one continuous campaign.
\systemname specified ACE-2, wrote its RTL, built its verification environment,
and drove it through mapped synthesis and static timing without a human author of
record for any of those artifacts. The result is an inference accelerator that
executes Qwen2.5-0.5B end to end in W4A8, certified against a fixed scope.

\paragraph{Functional closure.} The accepted design runs the full 24-layer
Qwen2.5-0.5B W4A8 command integration. Layer 0 matches the reference on all 18
ordered operators exactly, and the two-token runtime completes
13{,}914/13{,}914 commands over 1{,}240{,}410{,}384 simulator cycles with
generated token identifiers $[0,0]$ and no first failure recorded. The runtime
package, command log, and progress journal are each bound by content hash, so the
functional claim resolves to specific bytes rather than to a summary.

\paragraph{Physical closure.} Canonical SKY130 HD mapped synthesis and OpenSTA at
TT 25\,\textdegree C / 1.80\,V report 62{,}283 cells and 0.614\,mm$^2$ of
non-SRAM area against an operator-set cap of 2.0\,mm$^2$, with $+0.6966$\,ns
detailed setup slack, 0.00\,ns worst negative slack, 0.00\,ns total negative
slack, and a 10.000\,ns clock period. Both operator-owned targets---the area cap
and the 100\,MHz floor---pass without relaxation, and the passing packet is the
sole exactly-once canonical run rather than the best of several attempts.

\paragraph{What the gate produced.} The certification is dated 2026-08-04 and was
issued by a Reviewer that did not perform the work, bound to one RTL tree hash.
Each accepted RTL repair carries its own reviewer binding together with the
originating and resulting tree hashes, so the certified design is reachable from
the audit trail rather than asserted alongside it. The more informative artifact,
for the argument of this report, is what the certificate refuses to say. It
enumerates its own exclusions: no routed timing, no power signoff, no DRC/LVS, no
GDS or tapeout, no silicon validation, no generation beyond two tokens, no
external deployment interfaces, and no FPGA prototype. That list was produced by
the same gate that admitted the positive claims. An unbounded runtime would have
reported a chip; this one reported a chip and the exact perimeter of the evidence
supporting it.

ACE-2 therefore extends the pattern of Sections~\ref{sec:vertical-trace}
and~\ref{sec:paper-case-study} into a domain with an unusually unforgiving
verifier. The functional and physical claims are narrow by construction, and the
scope conditions are part of the output rather than a caveat added afterward.
